\section{An exact finite-family optimization}\label{sec:optimization}

For each completion script, the entropy bound gives many possible
rate pairs, one for each leaf law.  We compute their exact support
function by optimizing a weighted sum of rates.  A convex-geometric
step then selects the best point at a prescribed rectangular shape.

Fix a distinguished side \(L\) and a parameter \(t\ge0\).
For two different colors \(i,j\), let
\[
 \eta_{ij}=
 \begin{cases}
  1,&\text{if their owning sides include }L,\\
  t,&\text{otherwise}.
 \end{cases}
\]
A label read by those sides contributes with coefficient \(\eta_{ij}\)
to the score \(S+tT\).

For a conditional tree of color \(i\), denote its maximum score by
\(V_i\).
For base two, all three values start at zero.
For base three they start at
\begin{equation}\label{eq:initial-scores}
 V_{\colA}=\eta_{\colA\colC}\log2,\qquad
 V_{\colB}=\eta_{\colB\colC}\log2,\qquad
 V_{\colC}=0.
\end{equation}
A completion with center \(i\) and size \(m\) updates them simultaneously:
\begin{align}
 V_i^{\rm new}&=mV_i, \label{eq:score-center}\\
 V_j^{\rm new}
 &=\eta_{ij}\log\left[
   \left(e^{V_i/\eta_{ij}}+e^{V_j/\eta_{ij}}\right)^m
       -e^{mV_i/\eta_{ij}}\right],\quad j\ne i. \label{eq:score-other}
\end{align}
For a termination keeping \(i,j\), the score is
\begin{equation}\label{eq:termination-score}
 D_{ij}=\eta_{ij}\log
       \left(e^{V_i/\eta_{ij}}+e^{V_j/\eta_{ij}}\right).
\end{equation}
At weight zero, these expressions mean their continuous limits.
Equivalently, replace a weighted log-sum by the maximum score of the
allowed words.
In particular, in \eqref{eq:score-other} the all-\(i\) word remains
excluded; its zero-weight value is
\[
 \max_{1\le q\le m}\bigl((m-q)V_i+qV_j\bigr).
\]

\begin{definition}\label{def:support-function}
Let \(\calF_7\) contain both bases, every list of at most seven
completions with sizes \(2\) or \(3\) and arbitrary centers, every
termination pair, and every distinguished side \(L\).
Define
\begin{equation}\label{eq:F-definition}
 F(t)=\max_{\calF_7}\frac{D_{ij}(t)}{b\log R}.
\end{equation}
For an available square exponent estimate \(p\ge w(1)\), define
\[
 F_p(t)=\max\{F(t),(2+t)/p\}.
\]
\end{definition}

All ranges in this definition are finite.
The recurrences use only three conditional scores per script, even
though the underlying readable tree can be much larger.

The entropy maximization identity used below is the standard conjugacy
between log-sum-exp and negative entropy
\cite[Example~3.25]{boyd-vandenberghe}.  We include its elementary
proof and the conditional-tree argument needed for these recurrences.

\begin{lemma}\label{lem:score-optimality}
The values in \eqref{eq:initial-scores}--\eqref{eq:termination-score}
are the exact maxima of \(S+tT\) over their conditional or complete
leaf laws.  Thus \(F\) is the support function, in direction \((1,t)\),
of the downward convex hull of the normalized rate pairs
\((S/(b\log R),T/(b\log R))\) from \(\calF_7\).
\end{lemma}

\begin{proof}
For \(\lambda>0\), entropy optimization gives
\begin{equation}\label{eq:gibbs}
 \lambda\log\sum_d e^{a_d/\lambda}
  =\max_{(p_d)}
     \left\{\lambda\HH((p_d))+\sum_d p_d a_d\right\}.
\end{equation}
Jensen's inequality applied to the logarithm proves the upper bound;
the law proportional to \(e^{a_d/\lambda}\) attains it.

For a noncenter completion branch, sum over all length-\(m\)
words in \(\{i,j\}\) other than the all-\(i\) word.
For each such word the old optimal scores add across its slots.
Condition on the color word and, as each slot is processed, on the
histories already read in preceding slots.  The remaining law in
that slot is an admissible law for its old conditional tree, so its
score is at most the corresponding old optimum.  Summing gives the
upper bound even for dependent slot laws.  Independent conditional
maximizing laws attain it.
Equation \eqref{eq:gibbs} now gives \eqref{eq:score-other}.
The center and termination formulas follow in the same way.
The zero-weight cases follow by continuity on the compact leaf-law
simplex.

Taking the maximum over scripts gives the claimed support values;
convexification and downward closure do not change them for
nonnegative weights.
\end{proof}

The geometric step uses the supporting-halfspace description of a
closed convex set \cite[Theorem~8.24]{rockafellar-wets}.  The reduction
to a ray and the restriction of the parameter interval are proved
below for this family.  Let \(\calH_F\) be the closed downward convex
hull of the normalized rate pairs in Lemma~\ref{lem:score-optimality},
and let \(\calH_{F_p}\) also include the square rate point
\((2/p,1/p)\) before taking that hull.  These hulls are taken in the
nonnegative quadrant.  Their intersection with the ray
\((S/C,T/C)=(2d,kd)\) determines the best balanced bound supplied by
these rates.

\begin{theorem}\label{thm:shape-optimization}
For \(U=F\) or \(U=F_p\), and \(0\le k\le1\),
the largest feasible scale on this ray is
\begin{equation}\label{eq:radial-optimum}
 d_U(k):=\max\{d\ge0:(2d,kd)\in\calH_U\}
       =\min_{0\le t\le1}\frac{U(t)}{2+kt}.
\end{equation}
Consequently,
\begin{equation}\label{eq:functional-bound}
 w(k)\le\frac1{d_U(k)}
      =\max_{0\le t\le1}\frac{2+kt}{U(t)}.
\end{equation}
A maximizing parameter is the leftmost \(t\in[0,1]\) for which
\begin{equation}\label{eq:contact-rule}
 (2+kt)U'_+(t)\ge kU(t),
\end{equation}
or \(t=1\) if no such point exists.
\end{theorem}

\begin{proof}
By \cref{cor:balance}, a point \((2d,kd)\in\calH_U\) with \(d>0\) yields
\(w(k)\le1/d\).
Supporting hyperplanes in nonnegative directions give
\[
 d_U(k)
    =\inf_{t\ge0}\frac{U(t)}{2+kt}.
\]
The pure second-coordinate constraint is the limit as \(t\to\infty\).
One may equivalently extend the hull downward to all of \(\R^2\);
every finite supporting normal is then nonnegative.
Continuity of \(W\) justifies the closed hull and boundary points.

At \(t=1\), the score is total leaf entropy and is independent of
which side is distinguished.
Average a maximizing construction over the three choices of that side.
This gives the feasible point
\[
 (2U(1)/3,U(1)/3).
\]
The square point already has that ratio if it is a maximizer.
Consequently \(U(t)\ge(2+t)U(1)/3\), and for \(t\ge1\), \(k\le1\),
\[
 \frac{2+kt}{U(t)}
 \le \frac{3(2+kt)}{(2+t)U(1)}
 \le \frac{2+k}{U(1)}.
\]
The last value occurs at \(t=1\), so larger \(t\) can be omitted.
The averaged point is strictly positive in both coordinates, and
\(U\) is positive and continuous on \([0,1]\).
Thus the infimum is the minimum in \eqref{eq:radial-optimum}; taking
reciprocals proves \eqref{eq:functional-bound}.

Finally \(U\) is convex as a support function.
Put \(G(t)=(2+kt)U'_+(t)-kU(t)\).
For \(0\le x<y\), the convex secant inequality gives
\(U(y)-U(x)\le(y-x)U'_+(y)\), and hence
\[
 G(y)-G(x)
 \ge (2+kx)\bigl(U'_+(y)-U'_+(x)\bigr)\ge0.
\]
Thus the right derivative of \(U(t)/(2+kt)\), whose sign is the sign
of \(G(t)\), changes from negative to nonnegative at the first
contact in \eqref{eq:contact-rule}.  This identifies its minimum,
including endpoints, upward slope jumps and contact intervals.
\end{proof}

The theorem supplies a finite mathematical prescription for every
aspect ratio.
For an active script, \eqref{eq:gibbs} gives exposed leaf laws.
At a nonsmooth contact, convex combinations of points on the exposed
face meet the desired ray.  Coordinatewise decrease is also allowed.
These laws and combinations can be approximated as in
\cref{thm:entropy}.
No assertion that one center sequence is numerically optimal is needed.

Smaller script lists give valid subfamilies, while allowing more
completions gives further valid bounds.
The truncation to \(t\le1\) uses all three distinguished-side choices;
they should be retained when such subfamilies are formed.
